\documentclass[11pt,a4paper]{article}
\usepackage[utf8]{inputenc}
\usepackage[T1]{fontenc}
\usepackage[english]{babel}
\usepackage{lmodern}
\usepackage{geometry}
\usepackage{xcolor}
\usepackage{tocloft}
\usepackage{hyperref}
\usepackage{enumitem}
\usepackage{parskip}
\usepackage{booktabs}
\usepackage{longtable}
\usepackage{fancyhdr}
\usepackage{lastpage}
\usepackage{listings}
\usepackage{titlesec}
\geometry{margin=2.15cm}
\setlength{\headheight}{14pt}
\setlength{\emergencystretch}{3em}
\setlist{itemsep=3pt,parsep=1pt,topsep=4pt}
\definecolor{titleblue}{HTML}{123B5D}
\definecolor{softblue}{HTML}{EDF5FB}
\definecolor{softgray}{HTML}{F5F7FA}
\definecolor{accent}{HTML}{0F6E8C}
\hypersetup{
  colorlinks=true,linkcolor=titleblue,urlcolor=accent,
  pdftitle={Scientific Data Analysis v2.2: English publication edition},
  pdfauthor={interemi},
  pdfsubject={Plugin-eval evaluation hardening}
}
\titleformat{\section}{\Large\bfseries\color{titleblue}}{\thesection.}{0.6em}{}
\titleformat{\subsection}{\large\bfseries\color{accent}}{\thesubsection.}{0.6em}{}
\setlength{\cftsecnumwidth}{2.8em}
\tocloftpagestyle{fancy}
\lstdefinestyle{promptstyle}{
  basicstyle=\ttfamily\small,backgroundcolor=\color{softgray},
  frame=single,rulecolor=\color{softblue},breaklines=true,columns=fullflexible
}
\pagestyle{fancy}
\fancyhf{}
\fancyhead[L]{Scientific Data Analysis}
\fancyhead[R]{v2.2 guide: English edition}
\fancyfoot[C]{\thepage\ of \pageref{LastPage}}
\newcommand{\code}[1]{\texttt{\detokenize{#1}}}

\begin{document}
\pagenumbering{gobble}
\begin{titlepage}
\centering
\vspace*{1.0cm}
{\Huge\bfseries\color{titleblue} Evaluation Hardening Guide\\[0.15cm]
Scientific Data Analysis v2.2\par}
\vspace{0.65cm}
{\Large A compact, assessable backend compatible with Scientific Workbench\par}
\vspace{0.8cm}
{\large Maintenance: a clean active bundle, historical routing stub,
regression wrappers, and post-sync plugin-eval\par}
\vspace{0.8cm}
{\large English publication edition: September 23, 2026\par}
\vspace{0.5cm}
\begin{minipage}{0.93\textwidth}
\small\raggedright
Translation of the preserved private v2.2 guide at commit \code{83dd0a5}.
Original TeX SHA-256:
\nolinkurl{9f9dd5ab069c07b65ba905586de2937caf4f13d024f1c80eec58afe3fde62a95}.
Private source:
\nolinkurl{skills/scientific-data-maintainer/.cache/active_references/guides/guia_scientific_data_analysis_v2_2.tex}.

Recorded results and closure steps are historical, not current verification.
This release guide does not synchronize or modify Scientific Workbench;
the app remains outside its scope. Historical commands are retained for
traceability: use fresh output directories when reproducing them, without
overwriting earlier evidence. The original remains unchanged privately.
\end{minipage}
\vfill
{\large Prepared with Codex\par}
{\large Copyright 2026 interemi\par}
\end{titlepage}

\clearpage
\pagenumbering{roman}
\fancyfoot[C]{\thepage}
\begin{abstract}
\code{scientific-data-analysis} v2.2 hardens evaluation on top of the v2.1
compact router and integrated v2.0 baseline. Its purpose remains reproducible,
non-destructive, traceable scientific and technical workflows. Packaging and
evaluator visibility change: temporary files, caches, and substantial historical
routing are excluded from the active bundle; representative regressions become
more visible to generic tools; and the guide advances to v2.2.

The skill remains the deterministic local backend, with Scientific Workbench
as the compatible app. v2.2 neither edits nor synchronizes the app.
\end{abstract}
\tableofcontents
\newpage
\pagenumbering{arabic}
\fancyfoot[C]{\thepage\ of \pageref{LastPage}}

\section{Executive overview}
v2.2 adds no capability. It makes the installed skill easier to assess,
maintain, and load. The recorded release work:
\begin{itemize}[leftmargin=1.5em]
\item Keeps \code{SKILL.md} as a compact router.
\item Preserves the registry, JSON contracts, run bundles, artifact types, and public scripts.
\item Removes \code{tmp/}, \code{__pycache__/}, and historical outputs from the active bundle.
\item Replaces long v2.1 routing with a compatible stub.
\item Replaces internal \code{README.md} documents counted by plugin-eval as scattered documentation.
\item Archives historical notes without live references in \code{fixtures/archived_references/}.
\item Extends \code{unittest} wrappers for routes and app-facing failures.
\item Prepares a realistic, non-destructive \code{plugin-eval} benchmark.
\item Reruns \code{plugin-eval analyze} before and after synchronization.
\end{itemize}

This is v2.2, not v2.3 or v3.0. It changes neither scientific behavior nor the
app, and makes no claim of unverified capabilities. The historical cleanup
described here is not an instruction to delete preserved evidence today.

\section{From v2.1 to v2.2}
\begin{longtable}{p{0.20\linewidth}p{0.68\linewidth}}
\toprule
\textbf{Area} & \textbf{v2.2 change} \\
\midrule
Active bundle & Temporary files, caches, and audit residue stay outside the installed skill. \\
Historical routing & \code{references/v2-1-expanded-skill-routing.md} becomes a lightweight stub. \\
Internal documents & Non-primary \code{README.md} files become \code{README.txt}. \\
Historical archive & Notes not linked from live routes move to \code{fixtures/archived_references/}. \\
Visible regressions & \code{tests/test_audit_regression_wrappers.py} covers audits, public surface, general routes, FITS, and clean failures. \\
Plugin-eval & Records a baseline, measures deferred cost, and reassesses after synchronization. \\
Guide & Adds \code{guia_scientific_data_analysis_v2_2.tex} and \code{.pdf}. \\
\bottomrule
\end{longtable}

\section{What remains unchanged}
\begin{itemize}[leftmargin=1.5em]
\item No capability renaming or changes to \code{public_surface_registry.yaml}.
\item No changes to frozen artifact types or scientific script contracts.
\item Optional backends do not become core dependencies.
\item Scientific Workbench is neither edited nor synchronized.
\end{itemize}

\section{Evaluation rules}
Evaluate the installed skill with plugin-eval. These historical command
templates use a placeholder for the evaluator location; resolve that location
before use. The original account-specific home path is replaced with
\code{$HOME}.
\begin{lstlisting}[style=promptstyle]
node .../plugin-eval.js analyze "$HOME/.codex/skills/scientific-data-analysis" --format markdown
node .../plugin-eval.js explain-budget "$HOME/.codex/skills/scientific-data-analysis" --format markdown
\end{lstlisting}

Interpret the measurements accurately:
\begin{itemize}[leftmargin=1.5em]
\item \code{trigger_cost_tokens}: cost of activating the skill.
\item \code{invoke_cost_tokens}: cost of reading initial instructions.
\item \code{deferred_cost_tokens}: bundled material that should not always be loaded.
\item Python warnings are real technical debt, but do not all justify a large refactor during a packaging release.
\end{itemize}

\section{Historical v2.2 closure criteria}
\begin{longtable}{p{0.48\linewidth}p{0.40\linewidth}}
\toprule
\textbf{Gate} & \textbf{Criterion} \\
\midrule
\code{py_compile} & Changed scripts and wrappers compile. \\
\code{sync_public_surface_docs.py --check} & Registry and snapshot remain consistent. \\
\code{skill_surface_audit.py} & Public surface remains valid. \\
\code{unittest wrappers} & Representative regressions pass. \\
\code{portable_smoke_test.py --profile core} & Core smoke passes or blocks in a controlled, justified way. \\
\code{plugin-eval analyze} & Score improves or remaining warnings are documented. \\
\code{benchmark.json} & Concrete scenarios are ready; execution may be manual if cost/time is unreasonable. \\
\code{post-sync} & Installed validation matches editable validation. \\
\bottomrule
\end{longtable}

\section{Archive policy}
The original guide records preservation outside the skill root at the following
historical location, shown with a portable alias:
\begin{lstlisting}[style=promptstyle]
$PRIVATE_WORKSPACE/skill_archives/scientific-data-analysis-v2_2-evaluation-hardening-20260617_100731
\end{lstlisting}
\code{$PRIVATE_WORKSPACE} identifies the original private editable workspace.
This archive supports recovery and traceability; it is not a public contract.
Do not restore the entire archive into the installed skill.

\section{Normal usage after v2.2}
End-user usage remains unchanged:
\begin{itemize}[leftmargin=1.5em]
\item General tables: \code{profile_table.py}, \code{cross_domain_data_workbench.py}.
\item Document intake: \code{document_intake_workbench.py}.
\item Astronomy/FITS: \code{inspect_fits.py}, \code{fits_rgb_batch.py}, \code{photometry_noise_budget.py}.
\item Notebooks: \code{notebook_workbench.py}.
\item App-readiness: JSON envelopes, run bundles, and artifact types remain valid.
\end{itemize}

\section{Accepted debt}
The original release accepts warnings that do not justify a large v2.2 refactor:
\begin{itemize}[leftmargin=1.5em]
\item Some substantial scientific scripts still have high complexity.
\item Some long lines originate in contracts, tables, or app-facing messages.
\item Deep deduplication belongs to an architecture phase rather than an evaluation release.
\end{itemize}
Do not break correct scientific behavior merely to improve an automated rating.

\section{Recommended commands}
\begin{lstlisting}[style=promptstyle]
python3 -m py_compile scripts/*.py scripts/_internal/*.py tests/test_audit_regression_wrappers.py
python3 scripts/sync_public_surface_docs.py --check
python3 scripts/skill_surface_audit.py --skip-hygiene
python3 -m unittest tests.test_audit_regression_wrappers
python3 scripts/portable_smoke_test.py --profile core
\end{lstlisting}

\section{Conclusion}
v2.2 improves evaluation and maintenance while preserving the central purpose:
reproducible local-first analysis, original preservation, traceable outputs,
and compatibility with Scientific Workbench.
\end{document}
